(1) 对各质量元相对于质心的 \(\boldsymbol\xi\) 展开 \(-\mu m/|\boldsymbol R+\boldsymbol\xi|\)。一阶质心项为零，二阶姿态项是
\(\mu[{\rm tr}I-3I_{rr}]/(2R^3)\) 的等价形式（姿态无关的常数可减去）。对轴对称惯量张量 \(I_{rr}=I_\perp+(I_\parallel-I_\perp)\cos^2\theta\)，姿态相关势为
\[
\boxed{U_{gg}=\frac{3n^2}{2}(I_\parallel-I_\perp)\cos^2\theta}.
\]
由于 \(I_\parallel<I_\perp\)，径向 \(\theta=0\) 的势较切向 \(\pi/2\) 更低，能差 \(\Delta U=\frac32n^2(I_\perp-I_\parallel)\)。这里先减去轨道质心引力势而不对它求转矩。

(2) 惯性角速度为 \(n+\dot\theta\)，俯仰平面转动动能为 \(\frac12I_\perp(n+\dot\theta)^2\)，取
\(L=\frac12I_\perp(n+\dot\theta)^2-U_{gg}\)。因交叉项 \(I_\perp n\dot\theta\) 为全时间导数，其存在不改变Euler–Lagrange方程：
\[
I_\perp\ddot\theta=\frac{3n^2}{2}(I_\parallel-I_\perp)\sin2\theta .
\]
当 \(|\theta|\ll1\)，\(\ddot\theta+\omega_{\rm lib}^2\theta=0\)，\(\boxed{\omega_{\rm lib}^2=3n^2(I_\perp-I_\parallel)/I_\perp>0}\)，无须额外姿态控制即可线性中性稳定。

(3) 自治的共转系能量为 \(\mathcal J=\frac12I_\perp\dot\theta^2+U_{gg}(\theta)\)，随时间严格不变。径向势谷与切向鞍顶的差是 \(\Delta U\)。由 \(\frac12I_\perp\dot\theta_0^2=\Delta U\)，得到分界临界初角速率
\[
\boxed{|\dot\theta_0|_{\rm sep}=n\sqrt{3(I_\perp-I_\parallel)/I_\perp}}.
\]
低于此值在同一势谷内往复俯仰，高于此值可跨过切向顶点继续转动；等号对应趋近鞍点的无限周期分界轨道，不能称为有限时刻通过不稳定点。

(4) 若 \(I_\perp=I_\parallel\)，四极扭矩与小振动频率均为零，所有方向在本阶近似等价。若 \(n\to0\) 则恢复无重力梯度的自由取向，临界速率随之趋零。正俯仰小角时 \(I_\parallel-I_\perp<0\)，扭矩指向减小角度，核对了恢复力方向。无阻尼系统的振荡并不自动衰减至零角；要实现渐近稳定必须另引入耗散或主动控制。
